%%%%%%%%%%%%%%%%%%%%%%%%%%%%%%%%%%%%%%%%%%%%%%%%%%%%%
%%%%
%%%%   Calculates and sets the correct page
%%%%   numbers for a particular article.
%%%%
%%%%%%%%%%%%%%%%%%%%%%%%%%%%%%%%%%%%%%%%%%%%%%%%%%%%%
%%%%
%%%%   usage: Insert the following lines before
%%%%          the \begin{document} into the article
%%%%
%%%%            \input procnum
%%%%            \numbering{<database>}{<articleId>}
%%%%
%%%%   Replace <articleId> with the id for the current article
%%%%   and <database> by the name of the textfile containing
%%%%   an entry for each article:
%%%%
%%%%     \article[<optRemark>]{<paperId>}{noPages}{author}{title}
%%%%
%%%%%%%%%%%%%%%%%%%%%%%%%%%%%%%%%%%%%%%%%%%%%%%%%%%%%%%%

\newcounter{startPage}
\setcounter{startPage}{1}

%% Define a command for the entries in the database
%% of articles :
%%     A counter is created from the <articleId>
%%     and initialized to the current startPage
%%     The startPage counter is incremented by the
%%     size of the current article :
%%
\newcommand{\article}[5][]{
  \newcounter{#2}
  \setcounter{#2}{\value{startPage}}
  \addtocounter{startPage}{#3}
}


%% The argument to the following command will be
%% the <articleId>. The corresponding counter's
%% value tells where to start the numbering
%% of this article.

\newcommand{\numbering}[2]{
    \input{#1}
    \setcounter{page}{\value{#2}}
}

%%%%%%%%%%%%%%%%%%%%%%%%%%%%%%%%%%%%%%%%%%%%%%%%%%%%%
